% =============================================================================
% 2. fejezet — Szakirodalmi áttekintés
% =============================================================================

Ez a fejezet a bőrelváltozások gépi tanulással történő osztályozásának
szakirodalmi előzményeit tekinti át. A klinikai diagnosztikai gyakorlat és a
korai felismerés jelentőségének ismertetése után végigköveti a terület
módszertani fejlődését -- a kézi jellemzőkön alapuló klasszikus gépi tanulástól a
konvolúciós neurális hálókon át a Vision Transformer és a multimodális nagy
nyelvi modellek megjelenéséig --, majd bemutatja az ISIC adathalmazokat, amelyek
a dolgozat empirikus alapját adják.

% =============================================================================
\section{A bőrrák és a korai diagnózis jelentősége}
% =============================================================================

A bőrrák a leggyakrabban diagnosztizált rosszindulatú daganatok közé tartozik
világszerte. A Nemzetközi Rákkutató Ügynökség (IARC) GLOBOCAN-becslései szerint
2022-ben több mint 1,5 millió új bőrdaganatos eset fordult elő, ezen belül a
melanomák száma megközelítőleg 331~700 volt, mintegy 58~700 halálesettel
\cite{iarc2022melanoma}. Bár a melanoma a bőrrákos esetek kisebb hányadát teszi
ki, agresszív viselkedése és áttétképző hajlama miatt a bőrdaganatok okozta halálozás
túlnyomó részéért felel.

A betegség kimenetele döntően a felismerés stádiumától függ. A még lokalizált
stádiumban diagnosztizált melanoma ötéves relatív túlélési aránya közel 100\%,
míg távoli áttét megjelenése esetén ez mintegy 35\%-ra csökken
\cite{seer2024melanoma}. Ez a markáns különbség teszi a korai felismerést
népegészségügyi szempontból kiemelten fontossá: a diagnózis néhány hónapos
késedelme is jelentősen ronthatja a beteg túlélési esélyeit.

\subsection{A klinikai diagnosztikai folyamat}

A melanoma klinikai gyanúja jellemzően vizuális vizsgálat alapján merül fel,
amelyet dermatoszkópos értékelés egészít ki; a végleges diagnózist minden
esetben szövettani vizsgálat (hisztopatológia) erősíti meg \cite{garbe2022melanoma}.
A magas kockázatú betegek nyomon követésében a szekvenciális digitális
dermatoszkópia és a teljes testfelszín fényképezése segíti a korai melanomák
kiszűrését \cite{garbe2022melanoma}. Ez a folyamat magas szintű szakmai
tapasztalatot és speciális eszközöket igényel.

A pigmentált bőrelváltozások szabad szemmel történő, gyors megítélését a
klinikai gyakorlatban az úgynevezett ABCDE-szabály támogatja. A betűszó a
melanomára jellemző öt figyelmeztető jegyet foglalja össze: aszimmetria
(\textit{Asymmetry}), szabálytalan határvonal (\textit{Border}), színbeli
heterogenitás (\textit{Color}), 6~mm-t meghaladó átmérő (\textit{Diameter}),
valamint az elváltozás időbeli változása (\textit{Evolving}); az utóbbi, ötödik
kritériummal Abbasi és munkatársai egészítették ki az eredeti ABCD-szabályt,
hangsúlyozva a fejlődő elváltozások jelentőségét \cite{abbasi2004abcde}. Az
ABCDE-szabály egyszerűsége miatt mind a szakemberek, mind a betegek számára jól
használható, ugyanakkor önmagában nem helyettesíti a műszeres és szövettani
vizsgálatot. Az ABCDE-kritériumokat a \ref{fig:abcde}.~ábra szemlélteti.

\begin{figure}[H]
  \centering
  \includegraphics[width=\textwidth]{images/diagrams/abcde_chart.png}
  \caption{Az ABCDE-szabály öt kritériuma valós felvételeken: a felső sor
    jóindulatú (nem melanoma), az alsó sor melanoma eseteket szemléltet az
    aszimmetria (A), a határvonal (B), a szín (C), az átmérő (D) és az időbeli
    változás, evolúció (E) szempontjából. Forrás: Science Learning Hub
    (sciencelearn.org.nz).}
  \label{fig:abcde}
\end{figure}

\subsection{A hozzáférés korlátai}

A bőrgyógyászati szakellátáshoz való hozzáférés világszerte egyenetlen.
Vidéki térségekben és alacsonyabb jövedelmű országokban a szakorvosi
kapacitás korlátozott, ami késedelmes diagnózishoz vezet; de még fejlett
egészségügyi rendszerekben is jelentős várakozási idővel kell számolni a
szakvizsgálatig. Ez a hozzáférési szakadék motiválja az olyan automatizált,
akár okostelefonon is futó szűrőeszközök fejlesztését, amelyek elsődleges
triázsként jelezhetik, ha egy elváltozás szakorvosi vizsgálatot indokol. A
mobil egészségügyi (mHealth) megközelítés ígéretét korai vizsgálatok is
jelezték: Massone és munkatársai például a mobiltelefonnal készített képeken
alapuló melanoma-szűrés megvalósíthatóságát igazolták \cite{massone2007cellphone}.

% =============================================================================
\section{Klasszikus gépi tanulás a bőrelváltozások osztályozásában}
% =============================================================================

A mély tanulás elterjedése előtt a bőrelváltozás-osztályozó rendszerek jellemzően
kétlépcsős felépítést követtek: először kézzel tervezett jellemzőket
(\textit{handcrafted features}) nyertek ki a képből, majd ezeken egy klasszikus
gépi tanulási osztályozót -- például SVM-et, decision tree-t vagy
Random Forest-et -- tanítottak be. A jellemzők gyakran közvetlenül
az orvosi szempontokat tükrözték: az ABCDE-szabályhoz kapcsolódó alaki
(aszimmetria, határvonal), színbeli (szín-momentumok, hisztogramok) és textúra
jellemzőket (például a szürkeárnyalatos együttes-előfordulási mátrixból
(GLCM) számított mutatókat).

Kassem és munkatársai szisztematikus áttekintése a terület módszertani
fejlődését dokumentálja, és rámutat, hogy a klasszikus, kézi jellemzőkön alapuló
megközelítések teljesítményét erősen korlátozza a feature engineering minősége és a
felvételi körülmények (megvilágítás, felbontás, artefaktumok) iránti
érzékenység \cite{kassem2021review}. Ezek a rendszerek jellemzően jól
kontrollált, dermatoszkópos felvételeken működtek megbízhatóan, általánosító
képességük azonban korlátozott maradt. Mindazonáltal a kézi jellemzőkön alapuló
osztályozók a mai napig értékes \textit{viszonyítási alapként}
szolgálnak, mivel gyorsak, értelmezhetők, és megmutatják, mekkora teljesítmény
érhető el a reprezentáció automatikus tanulása nélkül.

% =============================================================================
\section{A mély tanulás megjelenése: konvolúciós neurális hálók}
% =============================================================================

A konvolúciós neurális hálók (CNN-ek) áttörést hoztak a képosztályozásban azzal,
hogy a jellemzőkinyerést és az osztályozást egyetlen, végponttól végpontig
tanítható modellbe integrálták, megszüntetve a kézi feature engineering
szükségességét. A bőrgyógyászati alkalmazások egyik úttörő munkájában Esteva és
munkatársai kimutatták, hogy egy nagy képadatbázison előtanított CNN a vizsgált
bőrrák-osztályozási feladatokban eléri a tapasztalt dermatológusok teljesítményét
\cite{esteva2017dermatologist}. Ezt követően több vizsgálat is megerősítette,
hogy a mély tanuláson alapuló rendszerek megközelíthetik, sőt bizonyos, jól
körülhatárolt feladatokban meg is haladhatják az emberi szakértők
teljesítményét: Haenssle és munkatársai egy CNN-t 58 dermatológussal vetettek
össze dermatoszkópos melanoma-felismerésben \cite{haenssle2018man}, Brinker és
munkatársai pedig arról számoltak be, hogy mély neurális hálók a melanoma
képi osztályozásában felülmúlták a bevont bőrgyógyászokat \cite{brinker2019deep}.

E rendszerek teljesítményének kulcsa a \textit{transzfertanulás} (transfer
learning): a modelleket először egy nagy, általános képadatbázison -- jellemzően
az ImageNet adathalmazon \cite{russakovsky2015imagenet} -- tanítják be, majd a
megtanult vizuális reprezentációkat a jóval kisebb orvosi adathalmazra finomhangolják
(fine-tuning). Az architektúrák fejlődése szintén hozzájárult az eredmények
javulásához: a reziduális kapcsolatokat bevezető ResNet \cite{he2016resnet}
lehetővé tette a nagyon mély hálók stabil tanítását, az EfficientNet pedig a
hálózat mélységének, szélességének és felbontásának együttes, kiegyensúlyozott
skálázásával ért el kedvező teljesítmény-paraméterszám arányt
\cite{tan2019efficientnet}.

A korai rendszerek fontos korlátja volt, hogy elsősorban dermatoszkópos, nagy
felbontású klinikai felvételeken működtek megbízhatóan, amelyek elkészítéséhez
speciális eszköz és szakértelem szükséges. Emellett a szelekciós torzítás is
jelentős problémát jelentett: az elérhető adathalmazok döntően már szakorvoshoz
került, klinikailag gyanús elváltozásokat tartalmaztak, ami tovább korlátozta a
modellek általánosíthatóságát hétköznapi felvételekre.

% =============================================================================
\section{A Vision Transformer megjelenése}
% =============================================================================

A természetesnyelv-feldolgozásban bevezetett Transformer-architektúra a
\textit{figyelmi} (attention) mechanizmusra épül, amely a bemenet elemei közötti
összes páronkénti kapcsolatot képes közvetlenül modellezni, függetlenül azok
távolságától \cite{vaswani2017attention}. Ezt az elvet a képosztályozásra a
Vision Transformer (ViT) ültette át: a képet rögzített méretű foltokra
(\textit{patch}) bontja, ezeket lineárisan beágyazza, majd egy Transformer-kódoló
dolgozza fel őket \cite{dosovitskiy2021vit}. A konvolúciós hálókkal szemben a ViT
nem rendelkezik beépített képi induktív torzítással (mint a lokalitás vagy az
eltolás-invariancia), így a globális összefüggéseket rugalmasabban tanulja meg,
cserébe viszont jellemzően nagyobb tanítóadatra vagy erőteljes előtanításra van
szüksége a jó általánosításhoz \cite{dosovitskiy2021vit}. Az architektúra
részletes elméleti tárgyalása a \ref{chpt:theory}.~fejezetben található.

A ViT-architektúra a bőrgyógyászati képosztályozásban is megjelent. Dagnaw és
munkatársai például Vision Transformer és konvolúciós modelleket vetettek össze
magyarázható mesterséges intelligencia (XAI) technikákkal kiegészítve, és a
ViT-alapú, illetve a CNN-alapú modellek esetében hasonló nagyságrendű pontosságot
mértek \cite{dagnaw2024vit}. Ez a megfigyelés alátámasztja, hogy a két
architektúracsalád összevetése -- pontosan kontrollált, azonos adat- és
tanítási feltételek mellett -- önmagában is érdemi tudományos kérdés, amelyre a
dolgozat kutatási része is kitér.

% =============================================================================
\section{Az ISIC adathalmaz}
% =============================================================================

A bőrelváltozás-osztályozás kutatását jelentősen előmozdította az annotált,
nyilvánosan elérhető adathalmazok megjelenése. A HAM10000 adathalmaz 10~015,
több forrásból származó dermatoszkópos felvételt tartalmaz a gyakori pigmentált
elváltozásokról, és egységes kiértékelési keretet biztosított a kutatóknak
\cite{tschandl2019ham10000}. Az International Skin Imaging Collaboration (ISIC)
versenysorozata tovább bővítette az elérhető képanyagot; Rotemberg és
munkatársai egy beteg-központú, a képek mellett klinikai metaadatokat is
tartalmazó adathalmazt mutattak be, hangsúlyozva a klinikai kontextus
szerepét a melanoma azonosításában \cite{rotemberg2021isic}.

A jelen dolgozat szempontjából kiemelt jelentőségű az ISIC 2024 verseny
(\textit{Skin Cancer Detection with 3D-TBP}), amely a teljes testfelszín
háromdimenziós fényképezéséből (3D total-body photography) kivágott, mintegy
400~000 elváltozás-képet tartalmaz, és binárisan -- rosszindulatú, illetve
jó\-indulatú vagy köztes kategóriába -- történő osztályozást tűz ki célul
\cite{isic2024}. Ezek a kivágott felvételek minőségükben közelebb állnak a
hétköznapi, akár okostelefonnal készíthető képekhez, mint a klasszikus,
nagy felbontású dermatoszkópos felvételek, ami a valós felhasználási
helyzethez közelíti a feladatot.

Az ISIC 2024 adathalmaz egyik meghatározó jellemzője a rendkívül erős
\textit{osztály-egyensúlytalanság}: a rosszindulatú esetek a teljes adathalmaznak
csupán töredékét (nagyságrendileg ezredrészét) teszik ki. Ennek megfelelően a
verseny nem a hagyományos pontosságot, hanem a 80\%-os érzékenység (true positive
rate) feletti tartományra korlátozott, részleges görbe alatti területet (partial
AUC) használja elsődleges metrikaként, amely orvosilag releváns módon a magas
érzékenységi tartományban méri a modellek teljesítményét \cite{isic2024}.
% MEGJEGYZÉS (Ervin): ide illeszthető néhány példakép az ISIC 2024 adathalmazból
% (jó- és rosszindulatú elváltozások). A versenyadat licencfeltételeit ellenőrizd
% a beillesztés előtt.
\begin{figure}[H]
  \centering
  \includegraphics[width=\textwidth]{images/literature/lesions.png}
  \caption{Reprezentatív minták az ISIC 2024 adathalmazból.}
  \label{fig:isic-samples}
\end{figure}

% =============================================================================
\section{Multimodális nagy nyelvi modellek a bőrgyógyászatban}
% =============================================================================

A legutóbbi években a multimodális nagy nyelvi modellek (multimodal large
language models, MLLM-ek) megjelenése új lehetőségeket nyitott az orvosi
képértelmezésben. Ezek a modellek -- mint a Google Gemini-családja
\cite{gemini2023} -- egyszerre képesek képi és szöveges bemenet feldolgozására,
és természetes nyelvű magyarázatot, indoklást is tudnak generálni.

A bőrgyógyászati alkalmazásukat több vizsgálat is értékelte. Cirone és
munkatársai a GPT-4~Vision és a LLaVA modelleket vetették össze melanoma és
jóindulatú elváltozások megkülönböztetésében, és azt találták, hogy a GPT-4V
lényegesen jobban teljesített (85\% pontosság a LLaVA 45\%-ával szemben),
ugyanakkor mindkét modell teljesítménye érzékelhetően romlott a sötétebb
bőrtónusok esetén -- ami a reprezentációs torzítás kockázatára és a sokszínűbb
adathalmazok szükségességére hívja fel a figyelmet \cite{cirone2024llm}. Zhou és
munkatársai célzottan bőrgyógyászati rendszert építettek (SkinGPT-4), amelyben
egy előtanított Vision Transformert egy nagy nyelvi modellel kapcsoltak össze,
és valós eseteken értékelték autonóm diagnosztikai képességét
\cite{zhou2024skingpt4}.

E munkák tanulsága kettős. Egyrészt a multimodális nyelvi modellek erőssége a
természetes nyelvű magyarázat és a felhasználóval folytatott párbeszéd, nem pedig
a megbízható, kalibrált osztályozás. Másrészt a teljesítményük populációk közötti
egyenetlensége óvatosságra int a klinikai felhasználásban. Ezek a megfontolások
közvetlenül indokolják a jelen dolgozatban követett tervezési döntést: a
\textit{klasszifikációt} egy célzottan betanított képi osztályozó modell (ViT)
végzi, míg a nagy nyelvi modell (Gemini) szerepe a kép minőségének ellenőrzésére,
az eredmény közérthető magyarázatára és a felhasználói kérdések megválaszolására
korlátozódik. E munkamegosztás részletes tárgyalása és empirikus alátámasztása a
dolgozat kutatási, illetve mérnöki részében található.
